\chapter{Modèles discrets et suites affines}
\section*{Exercice 6.1 -- Déplacement d'un équilibre}
Résoudre $u_{n+1}=0{,}75u_n+8$, $u_0=12$, et interpréter son équilibre.

\section*{Exercice 6.2 -- Sans équilibre}
Résoudre $u_{n+1}=u_n+5$, $u_0=7$. Pourquoi ne pas diviser par $1-a$ dans la formule affine générale ?

\section*{Exercice 6.3 -- Alternance convergente}
Résoudre $u_{n+1}=-u_n/2+3$, $u_0=0$. Décrire sa position par rapport à la limite.

\section*{Exercice 6.4 -- Remboursement fictif}
Une dette mathématique $D_0=1000$ évolue selon $D_{n+1}=1{,}01D_n-A$. Trouver la mensualité constante annulant $D_{12}$, sans arrondir les calculs intermédiaires.

\section*{Exercice 6.5 -- Forçage variable}
Résoudre $u_{n+1}=2u_n+n$, $u_0=0$, en vérifiant la formule $u_n=2^n-n-1$.

\section*{Exercice 6.6 -- Tours de Hanoï}
Si $H_n$ est le nombre minimal de déplacements de $n$ disques, justifier $H_{n+1}=2H_n+1$ et résoudre avec $H_0=0$.
